\section{Physics Gate：代理训练前的物理与离散验收}

本节定义生产代理训练之前必须通过的底层验证。其目的不是增加流程复杂度，而是保证神经网络学习的是一个已经物理闭合、数值收敛且定义一致的参数映射。任一关键 Gate 未通过时，不进入正式 surrogate training。

\subsection{Gate 0A：电磁 formulation 与开放域}
在代表性几何上验证 Maxwell formulation 与边界处理：
\begin{enumerate}
  \item 逐步扩大计算域或采用经过验证的开放边界/PML；
  \item 比较 $Z_{\rm field}$、互阻抗、体 Joule loss 与关键场量；
  \item 要求目标输出随边界位置/吸收层设置达到预设收敛阈值；
  \item 在低频 regime 下，将当前 formulation 与适当的低频 formulation（例如 $A$--$\phi$）做代表性对照，确认没有明显的低频病态或错误物理分支。
\end{enumerate}
只有通过该 Gate 后，才冻结 truth electromagnetic model。

\subsection{Gate 0B：端口定义、reciprocity、passivity 与功率闭合}
对每个 audit geometry 和多组复端口电流 $c$，检查：
\begin{align}
Z_{\rm field}^T&\approx Z_{\rm field},\\
\frac12\operatorname{Re}(c^HZc)
&\approx P_{\rm vol}+P_{\rm out}+P_{\rm wire},\\
\operatorname{Re}(c^HZc)&\ge -\varepsilon_{\rm passivity}.
\end{align}
测试电流必须覆盖单端口、同相、反相和 $\pm\pi/2$ 相位差。阻抗符号、source orientation、Joule Hodge 或 complex phase handling 任何一项错误都会在该 Gate 暴露。

除 truth audit 外，生产 surrogate 解码也必须按构造保证
\begin{equation}
\operatorname{Herm}(\widehat Z_{\rm field})\succeq0,
\end{equation}
不能只在有限 validation points 上检查后放任网络在域内其他点产生主动端口。

\subsection{Gate 0C：Joule 与 thermal projection 一致性}
从 Maxwell 使用的同一 $H_\sigma$ 构造空间体 Joule heat，并验证
\begin{equation}
\sum_c Q_{{\rm vol},c}
\approx\frac12E^HH_\sigma E.
\end{equation}
对每个 thermal mode 验证直接空间投影与 Hermitian tensor contraction 一致：
\begin{equation}
\phi_j^TQ_{\rm vol}
\approx\operatorname{Re}(c^HH_jc).
\end{equation}
禁止使用不能满足该能量闭合的独立 field reconstruction 定义 Joule heat。

\subsection{Gate 0D：几何合法性、source continuity 与参数连续性}
每个生产几何必须满足：
\begin{itemize}
  \item conductor 位于其设计包封/允许区域内；
  \item TX/RX conductor、package 等实体满足规定的无穿透/最小间隙约束；
  \item 与人工电磁/热边界保持足够裕量；
  \item material fraction 闭合且无非法重叠解释；
  \item source deposition 保持端口总电流归一、闭合回路拓扑与相应离散电流连续性。
\end{itemize}

对于闭合 filament coil，应显式验证离散 source 不产生非物理净电流端点；若离散体系具有 current-divergence/incidence operator，则相应闭合回路 source residual 应达到离散容差。仅保证 source 向量非零或积分幅值正确不足以证明 source 物理兼容。

此外，对几何参数做小扰动 $g\to g+\delta g$，检查 $B(g)$、material operators、$Z(g)$ 和 $H_j(g)$ 不出现由 cell midpoint assignment、离散采样点数量突变等造成的非物理跳跃。若存在 staircasing，应先改用连续、守恒的 line/source deposition 和足够平滑的 sub-cell quadrature，再训练代理。

\subsection{Gate 0E：空间离散收敛}
至少选取中心、边界和困难几何做 mesh refinement。分别监控
\begin{equation}
Z_{\rm field},\quad P_{\rm vol},\quad H_j,\quad T_{\max},\quad a_*.
\end{equation}
训练 truth 数据采用的网格必须位于工程目标所需的收敛区间。小 linear residual 不能替代 mesh convergence。

\subsection{Gate 0F：thermal 边界和 transient-aware ROM}
热模型边界条件必须有明确物理含义。若采用有限盒子上的 ambient Dirichlet 或 Robin boundary，应通过扩大 thermal domain、改变 boundary location 或与更高保真边界模型比较，确认长期温升对人工边界不过度敏感。

thermal basis 的自动 rank 必须覆盖 geometry、source pattern 与目标时间尺度。对 representative shifts $s_\ell\ge0$，直接检查在线 Galerkin resolvent residual
\begin{equation}
r_{\alpha\ell}
=Q_\alpha-(K+s_\ell M)\Phi
\left[\Phi^T(K+s_\ell M)\Phi\right]^{-1}\Phi^TQ_\alpha.
\end{equation}
shift 集必须包含 steady limit，并覆盖生产查询中的主要 inverse-time scales；只检查 $K^{-1}Q$ 与 $M^{-1}Q$ 两个端点不构成完整 transient 验证。

冻结 audit geometries 上还必须比较 full thermal 与 ROM 的 trajectory/steady-state outputs，包括至少一个短时、中间时间和接近稳态的时间尺度。若 wire heating 或其他热源幅值依赖温度，应在其生产状态范围内做 end-to-end trajectory audit。

\subsection{Gate 0G：模型假设边界}
当前线圈采用 filament current + 集中 AC resistance，而不是 full solid-conductor eddy-current model。因此必须明确其适用范围，并用更高保真或实验参考确认 skin/proximity effect 近似是否足够。

当前海水热模型若不含流体对流，只能代表所定义的 conduction/boundary heat-transfer model。若目标是长期真实海水温升，应单独验证自然/强制对流是否可以忽略或由有效边界/材料模型吸收。

当前 operating variable $c$ 表示规定的复端口电流。若真实系统由电压源、谐振补偿网络或负载闭环决定电流，则必须在 $Z(a,g)$ 外再接显式 circuit equations 求解 $c$；不能把任意固定电流结果直接解释为完整 WPT 电路工作点。

\subsection{Gate 通过后的数据冻结}
Physics Gate 通过后，应冻结并版本化：
\begin{itemize}
  \item PDE formulation 与相量/端口符号；
  \item 计算域、网格与边界条件；
  \item source/material deposition；
  \item constitutive parameters；
  \item Joule/impedance extraction；
  \item thermal full model、source family 与 ROM residual definition。
\end{itemize}
只有在上述版本固定后生成的 truth snapshots 才属于同一个可学习函数族。任何一项基础物理改变都应触发 truth dataset 版本更新，而不是直接在旧标签上继续训练。
